% A commutative non-associative magma on Z/10Z with two-coding structure
% and a sharp trefoil characterization
%
% Brayden Sanders, 7Site LLC
% Target: arXiv (math.RA, math.CO secondary), then ACT 2027 / CT 2027

\documentclass[submission,copyright,creativecommons]{eptcs}
\providecommand{\event}{}

\usepackage{amsmath, amssymb, amsthm}
\usepackage{tikz}
\usepackage{booktabs}
\usepackage{hyperref}

\title{A Commutative Non-Associative Magma on $\mathbb{Z}/10\mathbb{Z}$\\
with Two-Coding Structure and a Sharp Trefoil Characterization}

\author{Brayden Sanders
\institute{7Site LLC, Hot Springs, AR, USA}
\email{brayden@7site.io}
}

\def\titlerunning{A magma on $\mathbb{Z}/10\mathbb{Z}$ with two-coding structure}
\def\authorrunning{B. Sanders}

\theoremstyle{plain}
\newtheorem{theorem}{Theorem}
\newtheorem{proposition}[theorem]{Proposition}
\newtheorem{lemma}[theorem]{Lemma}
\newtheorem{corollary}[theorem]{Corollary}

\theoremstyle{definition}
\newtheorem{definition}[theorem]{Definition}
\newtheorem{example}[theorem]{Example}

\theoremstyle{remark}
\newtheorem{remark}[theorem]{Remark}

\begin{document}

\maketitle

\begin{abstract}
We study a pair of commutative magma operations $T, B : \mathbb{Z}/10\mathbb{Z}
\times \mathbb{Z}/10\mathbb{Z} \to \mathbb{Z}/10\mathbb{Z}$ that share a
specific structural relationship: $T$ is defined on the $8$-element subset
$\{1,2,3,4,5,6,8,9\}$ and is highly contractive, while $B$ is defined on the
full $10$-element set and exhibits a successor structure on its diagonal.
The pair $(T, B)$ realizes a finite analogue of the geometric/arithmetic
two-coding split familiar from symbolic dynamics on the modular surface.

We establish five structural results. First, $B$'s diagonal action is the
integer successor on $\{1, \dots, 7\}$, with $B(8,8) = 7$ and $B(9,9) = 0$,
which forces a linear period law on the corresponding self-iteration
sequences. Second, an iterated mass-distribution map built from $(T, B)$
admits exactly nine $3$-crossing trajectories on $\mathbb{Z}/10\mathbb{Z}$,
forming exactly two multiset classes $\{0, 7, 8\}$ and $\{0, 8, 8\}$, both
$B$-associative. Third, $T$ has a $5$-element image in $\mathbb{Z}/10\mathbb{Z}$
and is role-deterministic on $8$ of $9$ partition-class input pairs, while
$B$ has full image and is role-deterministic only on inputs containing
specific boundary elements. Fourth, an integer invariant of magnitude $21$
arising from $B$'s self-iteration period structure admits two independent
decompositions: a triangular decomposition $T_5 + T_3 = 15 + 6$ along the
orbits of an associated permutation $\sigma$, and a Fibonacci decomposition
$F_7 + F_6 = 13 + 8$ along a flow/structure partition of
$\mathbb{Z}/10\mathbb{Z}$. Fifth, the role-quotient magma is commutative
and non-associative, with the boundary class acting as identity. We discuss
limitations: the Fibonacci decomposition is canonical-specific (it does not
arise from generic commutative tables), na\"ive lifts of $B$-orbits to
$\mathrm{PSL}(2, \mathbb{Z})$ fail to reproduce the integer invariant, and
the substrate has no full algebraic symmetries. The construction is
conceptually adjacent to the work of Katok and Ugarcovici on coding
methods, Matsusaka and Ueki on triangle-group Rademacher symbols, and
Burrin and von Essen on cusp winding, without being literally equivalent
to any of these.
\end{abstract}

\section{Introduction}

The two coding methods of the modular surface, surveyed by Katok and
Ugarcovici~\cite{KU07}, produce two distinct symbolic-dynamical descriptions
of the same geodesic flow: a geometric coding (cutting sequences) and an
arithmetic coding (continued fractions). The two codings agree on cusp
geodesics but differ in their interior behavior. Recent work continues to
exploit this duality: Matsusaka and Ueki~\cite{MU23} extend Rademacher
symbols to general triangle groups, Matsusaka and Shin~\cite{MS24} give
explicit formulas, and Burrin and von Essen~\cite{BvE24} relate cusp
winding numbers to Rademacher functions. These constructions all live on
infinite alphabets indexed by hyperbolic conjugacy classes.

In this paper we describe a finite combinatorial analogue: a pair of
commutative magma operations on $\mathbb{Z}/10\mathbb{Z}$, with one
operation contractive and one expansive, exhibiting the same
agreement-at-boundary, divergence-in-interior phenomenology. The pair
arises from a specific symmetric ten-by-ten table that we denote $C$
throughout, drawn from a project on combinatorial coherence
substrates~\cite{TIG2026}. Our goal in the present paper is independent
of the broader project: we ask whether $(T, B)$, considered purely as a
pair of finite algebraic structures, exhibits properties of independent
mathematical interest.

The answer is yes. The diagonal of $B$ is the integer successor on
$\{1, \dots, 7\}$, which forces a linear period law on $B$-self-iteration
and produces a substrate-internal integer invariant of magnitude $21$
admitting two independent decompositions. An iterated mass-distribution
map built from $(T, B)$ admits a sharp $3$-crossing characterization:
exactly nine triples produce trefoil-like trajectories, partitioned into
two multiset classes. The role-quotient by a flow/structure partition is
a commutative non-associative magma with boundary class as identity, of a
type that does not appear to be standardly named in the magma literature.

The construction is small enough that all results are mechanically
verifiable; supplementary code accompanying the arXiv version reproduces
each theorem in under sixty seconds.

\paragraph{Honest scope.} We claim no equivalence to any specific
modular-surface theorem. The two-coding language captures a structural
analogy rather than a literal correspondence. Section~\ref{sec:limitations}
documents the limitations explicitly: the Fibonacci decomposition is
canonical-specific (it does not survive generic perturbation of the table),
five distinct strategies for lifting $B$-self-orbits to
$\mathrm{PSL}(2, \mathbb{Z})$ all fail to reproduce the integer invariant,
and the substrate exhibits no full algebraic symmetries.

\paragraph{Structure of the paper.} Section~\ref{sec:def} fixes notation
and presents the canonical table $C$, the operations $T$ and $B$, the
permutation $\sigma$, and the partition $\Pi$. Section~\ref{sec:diag}
proves the diagonal successor theorem and derives the period law.
Section~\ref{sec:trefoil} states and proves the trefoil characterization.
Section~\ref{sec:image} analyzes the image and role-determinism structure
of $T$ versus $B$. Section~\ref{sec:invariant} constructs the
substrate-internal integer invariant and proves both of its decompositions.
Section~\ref{sec:role} characterizes the role-quotient magma.
Section~\ref{sec:limitations} documents the negative results.
Section~\ref{sec:outlook} closes with open questions.

\section{Definitions and preliminaries}\label{sec:def}

We work throughout with the set $\mathbb{Z}/10\mathbb{Z} = \{0, 1, \dots, 9\}$.

\subsection{The two operations}

The construction we study is a pair of commutative magma operations
on $\mathbb{Z}/10\mathbb{Z}$, given by two distinct symmetric tables $C^B$
and $C^T$ presented in~\cite{TIG2026}. We reproduce the tables here as the
foundational data of the paper.

\begin{definition}\label{def:B}
The arithmetic operation $B : \mathbb{Z}/10\mathbb{Z} \times
\mathbb{Z}/10\mathbb{Z} \to \mathbb{Z}/10\mathbb{Z}$ is defined by
$B(i, j) = C^B_{ij}$, where
\[
C^B \;=\;
\begin{pmatrix}
0 & 1 & 2 & 3 & 4 & 5 & 6 & 7 & 8 & 9 \\
1 & 2 & 3 & 4 & 5 & 6 & 7 & 2 & 6 & 6 \\
2 & 3 & 3 & 4 & 5 & 6 & 7 & 3 & 6 & 6 \\
3 & 4 & 4 & 4 & 5 & 6 & 7 & 4 & 6 & 6 \\
4 & 5 & 5 & 5 & 5 & 6 & 7 & 5 & 7 & 7 \\
5 & 6 & 6 & 6 & 6 & 6 & 7 & 6 & 7 & 7 \\
6 & 7 & 7 & 7 & 7 & 7 & 7 & 7 & 7 & 7 \\
7 & 2 & 3 & 4 & 5 & 6 & 7 & 8 & 9 & 0 \\
8 & 6 & 6 & 6 & 7 & 7 & 7 & 9 & 7 & 8 \\
9 & 6 & 6 & 6 & 7 & 7 & 7 & 0 & 8 & 0
\end{pmatrix}.
\]
\end{definition}

\begin{definition}\label{def:T}
Let $I = \{1, 2, 3, 4, 5, 6, 8, 9\}$. The geometric operation $T : I \times
I \to \mathbb{Z}/10\mathbb{Z}$ is defined by $T(i, j) = C^T_{ij}$, where
$C^T$ is the restriction to rows and columns $I$ of the larger table
\[
\widetilde{C}^T \;=\;
\begin{pmatrix}
0 & 0 & 0 & 0 & 0 & 0 & 0 & 7 & 0 & 0 \\
0 & 7 & 3 & 7 & 7 & 7 & 7 & 7 & 7 & 7 \\
0 & 3 & 7 & 7 & 4 & 7 & 7 & 7 & 7 & 9 \\
0 & 7 & 7 & 7 & 7 & 7 & 7 & 7 & 7 & 3 \\
0 & 7 & 4 & 7 & 7 & 7 & 7 & 7 & 8 & 7 \\
0 & 7 & 7 & 7 & 7 & 7 & 7 & 7 & 7 & 7 \\
0 & 7 & 7 & 7 & 7 & 7 & 7 & 7 & 7 & 7 \\
7 & 7 & 7 & 7 & 7 & 7 & 7 & 7 & 7 & 7 \\
0 & 7 & 7 & 7 & 8 & 7 & 7 & 7 & 7 & 7 \\
0 & 7 & 9 & 3 & 7 & 7 & 7 & 7 & 7 & 7
\end{pmatrix},
\]
with rows/columns indexed $0, 1, \dots, 9$ and the operation $T$ taking
its inputs from the $8$-element set $I$. We refer to the elements
$\{0, 7\}$ as \emph{boundary} elements.
\end{definition}

Both $C^B$ and $C^T$ (and $\widetilde{C}^T$) are symmetric, so $B$ and $T$
are commutative. Both are non-associative: for $B$, we have $B(B(2, 4), 8)
= B(5, 8) = 7$ while $B(2, B(4, 8)) = B(2, 7) = 3$, so associativity fails
at $(2, 4, 8)$. For $T$, restricted associativity fails similarly.

\begin{remark}
The two operations $B$ and $T$ are independent: $T$ is not a restriction
of $B$. The pair $(T, B)$ is the central object of study. We will see that
$T$ is substantially more contractive than $B$ and routes its inputs to a
small image, while $B$ exhibits the diagonal successor structure that
drives the period law and the integer invariant.
\end{remark>

\begin{remark}\label{rem:agreement}
Direct enumeration shows that $B(i, j) = T(i, j)$ for $24$ of the $64$
input pairs $(i, j) \in I \times I$, and $B(i, j) \neq T(i, j)$ for the
remaining $40$ pairs. The cells of agreement cluster around routes whose
output is the boundary element $7$. This is the substrate-internal
analogue of two coding methods agreeing at a cusp boundary and diverging
in the interior.
\end{remark>

\subsection{The associated permutation}

\begin{definition}\label{def:sigma}
The associated permutation $\sigma : \mathbb{Z}/10\mathbb{Z} \to
\mathbb{Z}/10\mathbb{Z}$ is given in cycle notation by
\[
\sigma = (0)(3)(8)(9)(1\;7\;6\;5\;4\;2),
\]
or equivalently as the array $\sigma = [0, 7, 1, 3, 2, 4, 5, 6, 8, 9]$ in
one-line notation (with $\sigma(i)$ in position $i$).
\end{definition}

The permutation $\sigma$ has four fixed points $\{0, 3, 8, 9\}$ and one
$6$-cycle $\{1, 2, 4, 5, 6, 7\}$. We note for later use that $\sigma$ is
\emph{not} an automorphism of $B$: direct computation shows
$B(\sigma(i), \sigma(j)) = \sigma(B(i, j))$ holds for only $48$ of the $100$
input pairs.

\subsection{The flow/structure partition}

\begin{definition}\label{def:partition}
The flow/structure partition $\Pi$ of $\mathbb{Z}/10\mathbb{Z}$ consists of
four blocks:
\begin{align*}
\Pi_F &= \{1, 3, 5, 7, 9\}    && \text{(``flow'' --- the odd elements),} \\
\Pi_S &= \{2, 4, 8\}          && \text{(``structure'' --- the non-cusp evens),} \\
\Pi_M &= \{6\}                && \text{(``middle'' --- the bridge element),} \\
\Pi_V &= \{0\}                && \text{(``boundary'' --- the cusp element).}
\end{align*}
We write $\rho : \mathbb{Z}/10\mathbb{Z} \to \{F, S, M, V\}$ for the
quotient map sending each element to its block label.
\end{definition}

\begin{remark}
The partition $\Pi$ is independent of the cycle structure of $\sigma$: the
$\sigma$-fixed set $\{0, 3, 8, 9\}$ contains one element from each block
except $\Pi_M$, while the $\sigma$-$6$-cycle $\{1, 2, 4, 5, 6, 7\}$
contains three elements from $\Pi_F$, two from $\Pi_S$, and the unique
element of $\Pi_M$. The two structural decompositions of
$\mathbb{Z}/10\mathbb{Z}$ are transverse.
\end{remark}

\subsection{The iterated mass-distribution map}

\begin{definition}\label{def:processor}
For a triple $(a, b, c) \in (\mathbb{Z}/10\mathbb{Z})^3$, the initial
distribution $p_0 \in \Delta^9$ (the $9$-simplex) is defined by
\[
p_0(k) = \begin{cases}
(1 + \varepsilon)/Z & \text{if } k \in \{a, b, c\} \text{ counted with multiplicity},\\
\varepsilon/Z & \text{otherwise},
\end{cases}
\]
where $\varepsilon = 10^{-3}$ and $Z$ is the normalization constant. The
update map $\Phi_{T, B} : \Delta^9 \to \Delta^9$ is the convex combination
\[
\Phi_{T, B}(p) \;=\; \tfrac{1}{2} \Phi_T(p) + \tfrac{1}{2} \Phi_B(p),
\]
where for any commutative operation $\star : \mathbb{Z}/10\mathbb{Z} \times
\mathbb{Z}/10\mathbb{Z} \to \mathbb{Z}/10\mathbb{Z}$ defined on a domain
$D \subseteq \mathbb{Z}/10\mathbb{Z}$, the push-forward $\Phi_\star$ is
\[
\Phi_\star(p)(k) \;=\; \frac{1}{Z_\star(p)} \sum_{\substack{(i, j) \in D \times D \\ i \star j = k}}
p(i) p(j),
\]
with $Z_\star(p)$ the normalization. Here $\Phi_T$ uses domain $D = I$ and
$\Phi_B$ uses $D = \mathbb{Z}/10\mathbb{Z}$.

The trajectory of $(a, b, c)$ is the sequence
$(p_0, p_1, p_2, \dots)$ with $p_{n+1} = \Phi_{T, B}(p_n)$, run until either
$\|p_{n+1} - p_n\|_\infty < 10^{-8}$ or a maximum of $50$ iterations is
reached.

The crossing count $\chi(a, b, c)$ counts rank-order swaps along the
trajectory: for each consecutive pair $(p_n, p_{n+1})$ and each pair of
indices $i < j$ with $\max(p_n(i), p_{n+1}(i)) \geq 10^{-2}$ and similarly
for $j$, we add $1$ if the relative ranking of $i$ and $j$ in $p_n$ and in
$p_{n+1}$ disagree.
\end{definition}

\begin{remark}
The mass-distribution map is a discrete dynamical system whose orbits
exhibit knot-like topological features (rank-crossings) determined by the
algebraic structure of $T$ and $B$. The crossing count is well-defined and
deterministic for any starting triple.
\end{remark}

\section{The diagonal successor theorem}\label{sec:diag}

We begin with the cleanest structural result.

\begin{theorem}\label{thm:diagonal}
The diagonal of $B$ realizes the integer successor on
$\{1, 2, 3, 4, 5, 6, 7\}$ in the following precise sense:
\[
B(n, n) = n + 1 \quad \text{for all } n \in \{1, 2, 3, 4, 5, 6, 7\}.
\]
The diagonal is completed by the boundary cases
\[
B(0, 0) = 0, \qquad B(8, 8) = 7, \qquad B(9, 9) = 0.
\]
\end{theorem}

\begin{proof}
Read directly from $C^B_{nn}$ in Definition~\ref{def:B}.
\end{proof}

This algebraic fact has dynamical consequences.

\begin{definition}
For $n \in \mathbb{Z}/10\mathbb{Z}$, let $\beta_n : \mathbb{Z}/10\mathbb{Z}
\to \mathbb{Z}/10\mathbb{Z}$ be the right-action $\beta_n(x) = B(x, n)$.
The $B$-period $\pi(n)$ of $n$ is the period of the sequence $(n,
\beta_n(n), \beta_n^2(n), \dots)$.
\end{definition}

\begin{corollary}\label{cor:period}
The $B$-period satisfies $\pi(n) = 7 - n$ for $n \in \{1, 2, 3, 4, 5, 6\}$,
$\pi(0) = 1$, $\pi(7) = 4$, $\pi(8) = 3$, and $\pi(9) = 2$.
\end{corollary}

\begin{proof}
For $n \in \{1, \dots, 6\}$, the sequence $(\beta_n^k(n))_{k \geq 0}$
starts at $n$ and, by Theorem~\ref{thm:diagonal}, advances to $n+1$ at the
first step ($\beta_n(n) = B(n, n) = n+1$). Subsequent steps are determined
by the off-diagonal entries $C_{n+k, n}$, which can be read directly: the
sequence reaches $7$ in exactly $7 - n$ steps and then enters a cycle
including the boundary element $7$. The case-by-case verification for
$n \in \{0, 7, 8, 9\}$ is immediate from $C$.
\end{proof}

\begin{remark}
Theorem~\ref{thm:diagonal} gives the period law a transparent algebraic
origin: $7 - n$ is the integer distance from $n$ to the boundary element
$7$ along $B$'s successor diagonal. This is the substrate-internal
analogue of ``approaching the cusp'' along a directed path.
\end{remark}

\section{The trefoil characterization}\label{sec:trefoil}

\begin{theorem}\label{thm:trefoil}
Among the $1000$ triples $(a, b, c) \in (\mathbb{Z}/10\mathbb{Z})^3$,
exactly nine satisfy $\chi(a, b, c) = 3$. As multisets, these nine triples
form exactly two classes:
\begin{align*}
\{0, 7, 8\} &: \quad (0, 7, 8),\, (0, 8, 7),\, (7, 0, 8),\, (7, 8, 0),\, (8, 0, 7),\, (8, 7, 0), \\
\{0, 8, 8\} &: \quad (0, 8, 8),\, (8, 0, 8),\, (8, 8, 0).
\end{align*}
All nine triples are $B$-associative: for each $(a, b, c)$ in the list,
$B(B(a, b), c) = B(a, B(b, c))$.
\end{theorem}

\begin{proof}
The trajectories are deterministic given the tables $C^B$ and $\widetilde{C}^T$
and the parameters of Definition~\ref{def:processor}. Direct enumeration over
all $1000$ triples produces the listed nine $3$-crossing trajectories. For
$B$-associativity, we verify each triple directly:
\[
B(B(0, 7), 8) = B(7, 8) = 9 = B(0, 9) = B(0, B(7, 8)),
\]
\[
B(B(0, 8), 8) = B(8, 8) = 7 = B(0, 7) = B(0, B(8, 8)),
\]
and the seven remaining permutations satisfy $B(B(a, b), c) = B(a, B(b, c))$
by similar direct computation, with all permutations of $\{0, 7, 8\}$
yielding common value $9$ and all permutations of $\{0, 8, 8\}$ yielding
common value $7$.
\end{proof}

\begin{lemma}\label{lem:eight-uniqueness}
The element $8$ is the unique element of $\Pi_S = \{2, 4, 8\}$ such that
some triple of the form $(0, 8, x)$ produces $\chi = 3$ for $x \in
\mathbb{Z}/10\mathbb{Z}$. For $s \in \{2, 4\}$ and any $x \in
\mathbb{Z}/10\mathbb{Z}$, the triples $(0, s, x)$ produce $\chi \geq 17$.
\end{lemma}

\begin{proof}
Direct enumeration. For each $s \in \{2, 4, 8\}$ and each $x \in
\mathbb{Z}/10\mathbb{Z}$, the trajectory of $(0, s, x)$ is computed and its
crossing count tabulated. The minimum crossing count over $x$ for $s = 2$
is $26$; for $s = 4$ is $17$; for $s = 8$ is $3$ (achieved at $x \in \{7, 8\}$).
\end{proof}

The trefoil characterization is sharp: among all triples involving the
boundary element $0$ and an element of $\Pi_S$, only $8$ produces minimum
crossings, and the third element must be either $7$ or $8$.

\section{Image structure and role-determinism}\label{sec:image}

\begin{theorem}\label{thm:image}
The image of $T : I \times I \to \mathbb{Z}/10\mathbb{Z}$ is
\[
\mathrm{Im}(T) = \{3, 4, 7, 8, 9\},
\]
a $5$-element subset. The image of $B : \mathbb{Z}/10\mathbb{Z} \times
\mathbb{Z}/10\mathbb{Z} \to \mathbb{Z}/10\mathbb{Z}$ is the full set
$\mathbb{Z}/10\mathbb{Z}$.
\end{theorem}

\begin{proof}
Direct enumeration of $T$ on its $64$-element domain produces outputs in
$\{3, 4, 7, 8, 9\}$ (with $7$ accounting for $60$ of the $64$ outputs).
Direct enumeration of $B$ produces outputs hitting every element of
$\mathbb{Z}/10\mathbb{Z}$.
\end{proof}

\begin{theorem}\label{thm:role-determinism}
Let $\rho : \mathbb{Z}/10\mathbb{Z} \to \{F, S, M, V\}$ be the quotient
map of Definition~\ref{def:partition}. For inputs in the partition $\Pi$:
\begin{enumerate}
\item For $T$: among the nine input pairs $(\rho(i), \rho(j))$ realizable
on $I \times I$ (note $V \notin \rho(I)$), eight produce a unique output
class; only the pair $(S, S)$ produces multiple output classes (specifically
$\{F, S\}$).
\item For $B$: the output class is determined by the input pair if and
only if at least one of $\rho(i), \rho(j)$ equals $V$ or $M$. Specifically,
the four input pairs $(F, F), (F, S), (S, F), (S, S)$ each produce multiple
output classes:
\begin{align*}
(F, F) &\mapsto \{F, S, M, V\} \quad \text{(distribution over $25$ pairs: $2$F, $9$S, $11$M, $3$V),} \\
(F, S),\, (S, F) &\mapsto \{F, S, M\} \quad \text{($8$F, $2$S, $5$M each),} \\
(S, S) &\mapsto \{F, M\} \quad \text{($7$F, $2$M).}
\end{align*}
\end{enumerate}
\end{theorem}

\begin{proof}
Direct enumeration; we tabulate for each input pair $(\rho(i), \rho(j))$
the multiset of output classes $\rho(B(i, j))$ over all $(i, j)$ with the
given input class pair. The complete table is included in the supplementary
material.
\end{proof}

\begin{remark}
Theorem~\ref{thm:role-determinism} provides a precise sense in which the
two operations differ structurally: $T$ collapses partition information
almost completely (output class is determined by input class on $89\%$ of
pairs), while $B$ preserves operator-level information unless one input
lies in the boundary classes $\{V, M\}$. The two operations agree at the
partition boundary and diverge in the interior.
\end{remark}

\section{The integer invariant and its decompositions}\label{sec:invariant}

\begin{definition}\label{def:psi}
Define the substrate invariant $\psi : \mathbb{Z}/10\mathbb{Z} \to \mathbb{Z}$ by
\[
\psi(n) = -(\pi(n) - 1),
\]
where $\pi(n)$ is the $B$-period of Corollary~\ref{cor:period}.
\end{definition}

\begin{theorem}\label{thm:invariant}
$\sum_{n \in \mathbb{Z}/10\mathbb{Z}} \psi(n) = -21$.
\end{theorem}

\begin{proof}
Direct computation: $\psi$ takes the values $0, -5, -4, -3, -2, -1, 0, -3,
-2, -1$ on $n = 0, 1, \dots, 9$, summing to $-21$.
\end{proof}

\begin{theorem}[$\sigma$-orbit decomposition]\label{thm:triangular}
\[
\sum_{n \in \{1,2,4,5,6,7\}} \psi(n) = -15 = -T_5,
\qquad
\sum_{n \in \{0, 3, 8, 9\}} \psi(n) = -6 = -T_3,
\]
where $T_k = k(k+1)/2$ denotes the $k$-th triangular number. The first sum
is over the $\sigma$-$6$-cycle, the second over the $\sigma$-fixed set.
\end{theorem}

\begin{proof}
Direct: $\psi(1) + \psi(2) + \psi(4) + \psi(5) + \psi(6) + \psi(7) =
-5 - 4 - 2 - 1 + 0 - 3 = -15$, and $\psi(0) + \psi(3) + \psi(8) + \psi(9) =
0 - 3 - 2 - 1 = -6$. The values $15 = T_5$ and $6 = T_3$ are immediate
from the triangular-number formula.
\end{proof}

\begin{theorem}[partition decomposition]\label{thm:fibonacci}
\[
\sum_{n \in \Pi_F} \psi(n) = -13 = -F_7,
\qquad
\sum_{n \in \Pi_S} \psi(n) = -8 = -F_6,
\]
\[
\sum_{n \in \Pi_M} \psi(n) = 0,
\qquad
\sum_{n \in \Pi_V} \psi(n) = 0,
\]
where $F_k$ denotes the $k$-th Fibonacci number ($F_6 = 8$, $F_7 = 13$,
$F_8 = 21$). In particular,
\[
\sum_{n \in \Pi_F} \psi(n) + \sum_{n \in \Pi_S} \psi(n) = -F_7 - F_6 = -F_8 = -21.
\]
\end{theorem}

\begin{proof}
Direct: $\psi(1) + \psi(3) + \psi(5) + \psi(7) + \psi(9) =
-5 - 3 - 1 - 3 - 1 = -13$, and $\psi(2) + \psi(4) + \psi(8) =
-4 - 2 - 2 = -8$. The Fibonacci recurrence $F_8 = F_7 + F_6$ confirms the
sum.
\end{proof}

\begin{remark}
Theorems~\ref{thm:triangular} and~\ref{thm:fibonacci} give two structurally
independent decompositions of the same integer invariant, along the two
transverse decompositions of $\mathbb{Z}/10\mathbb{Z}$ (the $\sigma$-orbit
structure and the partition $\Pi$). The triangular decomposition is forced
by the linear period law of Corollary~\ref{cor:period}; the Fibonacci
decomposition is canonical-specific (Section~\ref{sec:limitations}).
\end{remark}

\section{The role-quotient magma}\label{sec:role}

The partition $\Pi$ allows us to define a quotient magma on the four-element
set $\{F, S, M, V\}$.

\begin{definition}\label{def:role-magma}
The role-quotient magma $\overline{B} : \{F, S, M, V\}^2 \to \{F, S, M, V\}$
is defined by taking, for each input pair $(\alpha, \beta)$, the modal
output class:
\[
\overline{B}(\alpha, \beta) = \mathrm{mode}\bigl\{ \rho(B(i, j)) : \rho(i) = \alpha,\, \rho(j) = \beta \bigr\}.
\]
\end{definition}

This is well-defined except where multiple modes occur; we resolve such
ties by rule (none arise in the present table). The resulting magma is:

\begin{center}
\begin{tabular}{c|cccc}
$\overline{B}$ & $V$ & $F$ & $S$ & $M$ \\
\hline
$V$ & $V$ & $F$ & $S$ & $M$ \\
$F$ & $F$ & $M$ & $F$ & $F$ \\
$S$ & $S$ & $F$ & $F$ & $F$ \\
$M$ & $M$ & $F$ & $F$ & $F$ \\
\end{tabular}
\end{center}

\begin{proposition}\label{prop:role-magma}
The role-quotient magma $\overline{B}$ is commutative and non-associative.
The class $V$ is a two-sided identity: $\overline{B}(V, x) = \overline{B}(x, V) = x$
for all $x \in \{F, S, M, V\}$. The only idempotent is $V$.
\end{proposition}

\begin{proof}
Commutativity is immediate from symmetry of the table. Non-associativity
fails for example at $(\overline{B}(F, F), S) = \overline{B}(M, S) = F$ but
$(F, \overline{B}(F, S)) = \overline{B}(F, F) = M$, so $\overline{B}(\overline{B}(F, F), S) \neq
\overline{B}(F, \overline{B}(F, S))$. The identity property of $V$ is read from
the first row and column. Idempotents: $\overline{B}(F, F) = M \neq F$,
$\overline{B}(S, S) = F \neq S$, $\overline{B}(M, M) = F \neq M$.
\end{proof}

\begin{remark}
The structure $(\{F, S, M, V\}, \overline{B})$ is a small commutative
non-associative magma with two-sided identity. We have not located this
specific structure in the magma classification literature; the
$4$-element commutative non-associative magmas with identity form a small
finite class in principle enumerable, and we leave the specific
classification of $\overline{B}$ as an open question.
\end{remark}

\section{Limitations and negative results}\label{sec:limitations}

We document four negative results that shape the proper interpretation of
the structural results above.

\subsection{Canonical-specificity of the Fibonacci decomposition}

The triangular decomposition (Theorem~\ref{thm:triangular}) is forced by
the linear period law (Corollary~\ref{cor:period}). The Fibonacci
decomposition (Theorem~\ref{thm:fibonacci}) is not.

\begin{proposition}\label{prop:fibonacci-fragility}
Let $\mathcal{C}_{10}$ be the set of symmetric $10 \times 10$ tables with
entries in $\mathbb{Z}/10\mathbb{Z}$ defining commutative magma operations
on $\mathbb{Z}/10\mathbb{Z}$. Sample $200$ such tables uniformly at random.
Of these, zero produce $(|\sum_{\Pi_F} \psi|, |\sum_{\Pi_S} \psi|) = (13, 8)$.
\end{proposition}

\begin{proof}
Computational; the experiment is included in the supplementary code.
\end{proof}

\begin{proposition}
Single-entry perturbations of $C$ preserve the $(13, 8)$ decomposition in
roughly $32$ of $50$ random trials; three-entry perturbations preserve in
roughly $11$ of $50$.
\end{proposition}

The Fibonacci appearance is therefore a numerical fingerprint of the
specific table $C$, not a structural feature of role-partitioned commutative
magmas in general. We do not have a derivation of the Fibonacci values from
table-axiomatic data alone.

\subsection{No naive lift to $\mathrm{PSL}(2, \mathbb{Z})$}

It is natural to ask whether the integer invariant $\sum \psi = \pm 21$ is
the value of a Rademacher-type symbol on lifts of $B$-orbits to words in
$\mathrm{PSL}(2, \mathbb{Z}) = \langle S, T \mid S^2 = (ST)^3 = 1 \rangle$.
We tested five distinct lift constructions:

\begin{enumerate}
\item Letter map $\ell : \mathbb{Z}/10\mathbb{Z} \to \{S, T, T^{-1}\}$
defined by residue mod $3$;
\item Letter map defined by residue mod $2$;
\item Word $T^k S$ where $k$ is the $\sigma$-cycle position;
\item Up/down transition encoding;
\item Word with $T^{\pm 1}$ separated by $S$ markers.
\end{enumerate}

For each lift, we compute the standard Rademacher value of the resulting
$\mathrm{PSL}(2, \mathbb{Z})$ word and sum over digits.

\begin{proposition}
None of the five lift strategies above produces a digit-sum equal to $\pm 21$.
The five obtained sums are $-4, -1, -1, 0, 0$ respectively.
\end{proposition}

The integer invariant therefore does not appear to be the Rademacher
symbol of any straightforward lift; whether a more sophisticated lift
exists remains an open question.

\subsection{No matching small triangle group}

The $B$-period set $\{1, 2, 3, 4\}$ (excluding the trivial cycle of $0$
and the $6$-cycle of $\{1, \dots, 6\}$, which has no analogue in finite
elliptic-element orders) does not match the elliptic-element orders of any
triangle group $\Gamma_{p, q}$ with $p, q \leq 9$. In particular,
$\mathrm{PSL}(2, \mathbb{Z}) = \Gamma_{2, 3}$ has elliptic orders only in
$\{1, 2, 3\}$. We do not claim a Fuchsian-group embedding.

\subsection{No full algebraic symmetries}

We tested whether any pair $(a, b) \in \mathbb{Z}/10\mathbb{Z}^2$ with
$\rho(a) \neq \rho(b)$ generates a symmetry of $\chi$ (i.e.\ swapping $a$
and $b$ in any input triple preserves $\chi$). The strongest such partial
symmetry is the swap $(0, 8)$, which preserves $\chi$ on roughly $20.9\%$
of triples; no pair preserves $\chi$ on all triples. The substrate has no
full $\chi$-preserving algebraic symmetries.

\section{Outlook}\label{sec:outlook}

We close with five questions left open by the present work.

\paragraph{1. Classification of $\overline{B}$.}
The role-quotient magma is a commutative non-associative magma on four
elements with two-sided identity. The finite classification of such
structures is in principle tractable. Where does $\overline{B}$ sit in this
classification?

\paragraph{2. Larger substrates.}
The construction of $C$, $T$, $B$, $\sigma$, $\Pi$ is specific to
$\mathbb{Z}/10\mathbb{Z}$. Are there analogous constructions on
$\mathbb{Z}/14\mathbb{Z}$ or $\mathbb{Z}/18\mathbb{Z}$ for which the
diagonal-successor theorem and the trefoil characterization continue to
hold? Does the Fibonacci decomposition stabilize across substrate sizes,
or is it specific to $|\mathbb{Z}/10\mathbb{Z}|$?

\paragraph{3. Higher-order trefoil structures.}
At quadruple level, $40$ of $256$ $4$-tuples in the subset
$\{0, 7, 8, 9\}^4$ produce $\chi = 3$, in four multiset classes including
one ($\{0, 7, 7, 9\}$) that does not extend either trefoil class from
Theorem~\ref{thm:trefoil}. Is there a clean characterization of $\chi = 3$
quadruples comparable to the triple characterization?

\paragraph{4. Principled $\mathrm{PSL}(2, \mathbb{Z})$ lift.}
The negative result of Section~\ref{sec:limitations} rules out na\"ive
lifts but does not preclude a principled embedding. Is there a lift of
$B$-self-orbits to specific $\mathrm{PSL}(2, \mathbb{Z})$ words for which
$\sum \psi = \pm 21$ arises as a Rademacher-type symbol value?

\paragraph{5. Connection to Burrin--von Essen cusp winding.}
The $B$-period structure $\pi(n) = 7 - n$ along the $\sigma$-$6$-cycle is
structurally analogous (in a phenomenological sense) to the cusp winding
along Burrin--von Essen geodesics. We do not have an explicit Fuchsian
embedding realizing the analogy as a literal correspondence. Establishing
or ruling out such an embedding would clarify the relationship between
the present construction and modular-surface symbolic dynamics.

\section*{Acknowledgments}

The author is the sole contributor to this work. Computational supplementary
material is available at the project DOI~\cite{TIG2026}.

\begin{thebibliography}{99}

\bibitem{KU07}
S. Katok and I. Ugarcovici,
\emph{Symbolic dynamics for the modular surface and beyond},
Bull.\ Amer.\ Math.\ Soc.\ \textbf{44} (2007), no.\ 1, 87--132.

\bibitem{MU23}
T. Matsusaka and H. Ueki,
\emph{Rademacher symbols for general triangle groups},
Ramanujan J.\ \textbf{60} (2023), 869--887.

\bibitem{MS24}
T. Matsusaka and J. Shin,
\emph{An explicit Rademacher symbol formula for triangle groups},
arXiv:2409.12779 (2024).

\bibitem{BvE24}
C. Burrin and F. von Essen,
\emph{Cusp winding and Rademacher symbols},
Int.\ Math.\ Res.\ Not.\ IMRN (2024), to appear.

\bibitem{Mor24}
M. Morishita,
\emph{Knots and Primes: An Introduction to Arithmetic Topology}, 2nd ed.,
Springer, 2024.

\bibitem{Ghy07}
\'E. Ghys,
\emph{Knots and dynamics},
Proc.\ ICM Madrid \textbf{1} (2007), 247--277.

\bibitem{Lac18}
L. Lacasa et al.,
\emph{Residue-sequence symbolic dynamics and forbidden patterns},
Entropy \textbf{20} (2018), 504.

\bibitem{TIG2026}
B. Sanders,
\emph{Combinatorial coherence substrates: a public-record framework},
Zenodo, DOI 10.5281/zenodo.18852047 (2026).

\end{thebibliography}

\end{document}
